\todo{Lecture 18}
\begin{example}[Explication du mouvement de la canette]
Le mouvement du centre de masse est décrit par 
$$
m\ddot{\boldsymbol{x}}_{\mathrm{CM}}=\sum _{i} \boldsymbol{F}^{\mathrm{ext}}_{i}
$$
et la rotation autour du centre de masse
$$
\frac{d\boldsymbol{L}}{dt}=\sum _{i}(\boldsymbol{x}_{i}-\boldsymbol{x}_{\mathrm{CM}})\times \boldsymbol{F}^{\mathrm{ext}}_{i}=\boldsymbol{\tau}_{\mathrm{CM}}^{\mathrm{ext}}
$$
si le corps est rigide
$$
I_{\mathrm{CM}} \dot{\boldsymbol{\omega}}=\boldsymbol{\tau}_{\mathrm{CM}}^{\mathrm{ext}}
$$

Ici, selon $x$ :
\begin{align*}
m \ddot{x}_{\mathrm{CM}} & =F_{1,x}+F_{2,x}+F_{f} \\
 I\dot{\omega} &R\cdot F_{f} \\
 \omega R & =\dot{x}_{\mathrm{CM}}
\end{align*}
ou $F_{1}$ est la force entrainée par la tige chargée.
\end{example}

\subsection{Effet de pointe}
Autour des pointes d'un conducteur, le champ $\boldsymbol{E}$ peut être très élevé.
\subsection{Traitement générale, equation de Laplace et de Poisson}
En principe, la loi de Gauss $\boldsymbol{\nabla E}=\frac{\rho _{\mathrm{el}}}{\varepsilon_{0}}$, permet de calculer $\boldsymbol{E}(\boldsymbol{r})$ a partir de $\rho _{\mathrm{el}}(\boldsymbol{r})$.
\begin{problem}
$\rho _{\mathrm{el}}(\boldsymbol{r})$ n'est pas connu partout.
\end{problem}
\begin{reminder}
$\boldsymbol{\nabla \nabla}=\Delta$.
\end{reminder}
\begin{example}
Soient trois conducteurs a des potentielles différentes. On connait le potentiel a la surface de chaque conducteur, mais on ne connait pas $\sigma _{\mathrm{el}}$. Si $\rho _{\mathrm{el}}=0$ dans l'espace a l'exterieur des conducteurs, alors $\boldsymbol{\nabla E}=0$, avec $\boldsymbol{E}=-\boldsymbol{\nabla}\phi$. Donc $\Delta \phi=0$, qui est l'equation de Laplace.

Si $\rho _{\mathrm{el}}\neq 0$ entres les conducteurs alors $\Delta \phi=-\frac{\rho _{\mathrm{el}}}{\varepsilon_{0}}$ est l'equation de Poisson.
\end{example}

On peut resoudre ce probleme en utilisant l'equation de Laplace (ou de Poisson) en dehors des conducteurs plus les conditions aux limites (valeurs de $\phi$ sur les conducteurs et a l'infini si le domaine est non-borne). 

Pour les conditions aux limites $\phi=\phi _{i}$ sur la surface du conducteur $i$ et $\phi \to0$ pour $\lVert \boldsymbol{r} \rVert\to \infty$ (si le domaine est non-borne), la solution de Poisson (ou Laplace) est unique.
\begin{example}[Cage de Faraday]
Une cage de Faraday est forme par un conducteur creux au potentiel $\phi_{1}$. A l'interier du conducteur $\Delta \phi=0$ et $\phi(\boldsymbol{r})=\phi_{1}$ est constante donc $\phi_{1}$ est l'unique solution a $\Delta \phi=0$. Alors $\boldsymbol{E}=-\boldsymbol{\nabla}\phi_{1}=0$ a l'interieur et $\sigma _{\mathrm{el}}=0$ sur la surface intérieure car $\lVert \boldsymbol{E} \rVert=\frac{\sigma _{\mathrm{el}}}{\varepsilon_{0}}$ a la surface d'un conducteur en électrostatique.
\end{example}